\section{Теорема об инвариантности ранга при элементарных преобразованиях. Нахождение ранга матрицы методом элементарных преобразований}

\begin{theorem}
При элементарных преобразованиях строк или столбцов ранг матрицы не изменяется.
\end{theorem}

\begin{proof}
Достаточно доказать утверждение для преобразований строк. Пусть $A \to \tilde{A}$.
\begin{enumerate}
    \item \textit{Перестановка строк.} Очевидно, что множество миноров и их значения не меняются, следовательно, $\operatorname{rang} A = \operatorname{rang} \tilde{A}$.
    \item \textit{Умножение $i$-й строки на $\lambda \neq 0$.} 
    Если базисный минор $M$ не содержит $i$-ю строку, то $M$ не меняется. Если $M$ содержит $i$-ю строку, то новый минор $\tilde{M}$ получается из $M$ умножением одной строки на $\lambda$. По свойству определителей $\tilde{M} = \lambda M$. Так как $\lambda \neq 0$ и $M \neq 0$, то $\tilde{M} \neq 0$. Ранг сохраняется.
    \item \textit{Прибавление к $i$-й строке $j$-й строки, умноженной на $\lambda$.} 
    Если базисный минор $M$ не содержит $i$-ю строку, он не меняется. Если содержит $i$-ю строку, то по свойству определителей прибавление к строке другой строки, умноженной на число, не меняет значение определителя: $\tilde{M} = M \neq 0$. Если же минор содержит обе строки, он также не меняется. Следовательно, максимальный порядок ненулевого минора (ранг) остается прежним.
\end{enumerate}
\hfill$\blacksquare$
\end{proof}

\begin{definition}
Матрица называется \textbf{ступенчатой (трапецевидной)}, если:
\begin{enumerate}
    \item Все нулевые строки (если они есть) расположены в самом низу матрицы.
    \item Ведущий (первый ненулевой слева) элемент каждой ненулевой строки находится строго правее ведущего элемента предыдущей строки.
\end{enumerate}
В общем виде такая матрица размера $m \times n$ ранга $r$ выглядит следующим образом:
$$
A_{\text{ступ}} = \begin{pmatrix}
a_{11} & a_{12} & \dots & a_{1r} & \dots & a_{1n} \\
0 & a_{22} & \dots & a_{2r} & \dots & a_{2n} \\
\vdots & \vdots & \ddots & \vdots & & \vdots \\
0 & 0 & \dots & a_{rr} & \dots & a_{rn} \\
0 & 0 & \dots & 0 & \dots & 0 \\
\vdots & \vdots & & \vdots & & \vdots \\
0 & 0 & \dots & 0 & \dots & 0
\end{pmatrix}, \quad \text{где } a_{11}, a_{22}, \dots, a_{rr} \neq 0.
$$
\end{definition}

\begin{theorem}
Всякую ненулевую прямоугольную матрицу путем элементарных преобразований строк и перестановок столбцов можно привести к ступенчатому виду.
\end{theorem}

\begin{proof}
Алгоритм приведения (метод Гаусса):
\begin{enumerate}
    \item Если $a_{11} = 0$, то перестановкой строк и/или столбцов добиваемся, чтобы на месте $(1,1)$ стоял ненулевой элемент.
    \item Умножаем первую строку на $1/a_{11}$, чтобы сделать ведущий элемент равным $1$ (этот шаг опционален для нахождения ранга, но полезен для приведения к каноническому виду).
    \item Прибавляем к каждой $i$-й строке ($i > 1$) первую строку, умноженную на $-a_{i1}/a_{11}$. В результате все элементы первого столбца ниже $a_{11}$ станут равны нулю.
    \item Применяем тот же процесс к подматрице, полученной вычеркиванием первой строки и первого столбца, и так далее, пока не исчерпаем все строки или столбцы.
\end{enumerate}
В результате получим матрицу ступенчатого вида, ранг которой равен количеству ненулевых строк.
\hfill$\blacksquare$
\end{proof}

\begin{remark}
Для нахождения ранга матрицы достаточно привести ее к ступенчатому виду с помощью элементарных преобразований строк. Ранг будет равен числу ненулевых строк в полученной матрице.
\end{remark}
